\documentclass[conference]{IEEEtran}
\usepackage{cite}
\usepackage{amsmath,amssymb,amsfonts}
\usepackage{algorithmic}
\usepackage{graphicx}
\usepackage{textcomp}
\usepackage{xcolor}
\usepackage{float}
\usepackage{hyperref}
\usepackage{booktabs}
\usepackage{listings}
\usepackage{tikz}
\usetikzlibrary{shapes,arrows,positioning}

% --- Custom Configurations ---
\makeatletter
\newcommand{\linebreakand}{%
  \end{@IEEEauthorhalign}
  \hfill\mbox{}\par
  \mbox{}\hfill\begin{@IEEEauthorhalign}
}
\makeatother

% Image placeholder generator
\newcommand{\imageplaceholder}[2]{
    \fbox{\begin{minipage}{#1}
        \centering
        \vspace{#2}
        \textbf{Visual Representation Placeholder}\\
        \small{(Upload target file to Overleaf workspace)}
        \vspace{#2}
    \end{minipage}}
}

\begin{document}

\title{Multi-Tiered Object Detection and Optical Character Recognition Ensemble with Combinatorial Error Correction for Automatic Wagon Parameter Auditing in Indian Railways}

\author{
    \IEEEauthorblockN{Dr. Vipul Kumar Mishra}
    \IEEEauthorblockA{Associate Professor, School of Technology,\\
    Gati Shakti Vishwavidyalaya,\\
    Vadodara, India - 390004\\
    Email: vipul.mishra@gsv.ac.in}
    \and
    \IEEEauthorblockN{Gagan Jain}
    \IEEEauthorblockA{UG Student, School of Technology,\\
    Gati Shakti Vishwavidyalaya,\\
    Vadodara, India - 390004\\
    Email: gagan.jain\textunderscore btech23@gsv.ac.in}
    \and
    \IEEEauthorblockN{Dev Yadav}
    \IEEEauthorblockA{UG Student, School of Technology,\\
    Gati Shakti Vishwavidyalaya,\\
    Vadodara, India - 390004\\
    Email: dev.yadav\textunderscore btech23@gsv.ac.in}
    \linebreakand 
    \IEEEauthorblockN{Chakshu Vashisth}
    \IEEEauthorblockA{UG Student, School of Technology,\\
    Gati Shakti Vishwavidyalaya,\\
    Vadodara, India - 390004\\
    Email: chakshu.vashisth\textunderscore btech23@gsv.ac.in}
    \and
    \IEEEauthorblockN{Ankit Gupta}
    \IEEEauthorblockA{UG Student, School of Technology,\\
    Gati Shakti Vishwavidyalaya,\\
    Vadodara, India - 390004\\
    Email: ankit.gupta\textunderscore btech23@gsv.ac.in}
}

\maketitle

\begin{abstract}
Efficient auditing of rolling stock parameters and maintenance safety stencils is a critical logistics challenge for heavy haul railway networks. This paper presents a deep learning based automation pipeline designed for Indian Railways (CRIS) to localize and parse wagon parameters and audit maintenance stencils from high-resolution images. The system employs a custom-trained YOLOv11m object detection model to localize five parameter classes: wagon number, wagon type, owner name, operating zone, and tare load details. A secondary YOLOv11m classifier audits the presence of scheduled safety maintenance stencils. To process degraded text on metal wagon bodies caused by paint peeling, rust, and dirt, we introduce a 5-tier fallback OCR ensemble strategy combining preprocessed PaddleOCR, raw PaddleOCR, EasyOCR, and whitelisted Tesseract. The localized outputs are cross-validated using the Indian Railways standard S4 Mod-10 checksum algorithm. When characters are degraded, a combinatorial search algorithm utilizes Levenshtein distance matrix matching against the master ROAMS registry to reconstruct the most probable wagon number. The system is deployed locally via a FastAPI server, an isolated PaddleOCR GPU microservice, and a React web dashboard. Experimental evaluations on real-world rolling stock demonstrate high parsing accuracy under extreme environmental noise.
\end{abstract}

\begin{IEEEkeywords}
Computer Vision, Object Detection, Optical Character Recognition, YOLOv11m, S4 Mod-10 Checksum, Indian Railways, Intelligent Infrastructure.
\end{IEEEkeywords}

\section{Introduction}
\label{sec:introduction}
The Indian Railways freight network is among the largest in the world, managing hundreds of thousands of active wagons. Keeping track of rolling stock parameters (such as the 11-digit wagon number, wagon type, owner name, operating zone, and tare load details) is crucial for asset tracking, scheduling, and preventive maintenance. Historically, wagon tracking has relied on manual visual inspections at yards and stations. Station staff or train examiners write down these parameters and enter them manually into central databases. This manual audit workflow is slow and prone to human errors, causing mismatches between the physical wagon state and the database registries.

Furthermore, wagons operate continuously in harsh environments. They are exposed to intense sunlight, rain, coal dust, industrial pollution, and physical impact. Consequently, the painted markings on their steel bodies degrade over time. Faded text, severe rust streaks, dirt layers, and high-contrast shadows cross-cutting the text boxes present major challenges for standard, single-stage optical character recognition (OCR) models.

To address these challenges, we have developed a robust, automated end-to-end deep learning pipeline. The core contributions of this work are:
\begin{enumerate}
    \item A two-stage localization strategy using custom-trained YOLOv11m object detectors to extract the parameter text regions and audit the presence of safety maintenance stencils.
    \item A robust 5-tier fallback OCR ensemble that sequentially processes text crops using custom image enhancements, character split-reduction models, and whitelisted layouts.
    \item A mathematical error-correction module that combines the S4 Mod-10 checksum algorithm with Levenshtein distance prefix lookup against the master ROAMS database to repair degraded digits.
    \item A multi-service deployment layout consisting of a FastAPI server, a React web interface, and an isolated GPU text recognition service optimized to prevent oneDNN thread crashes in Windows environments.
\end{enumerate}

\section{System Architecture and Bounding Box Localization}
\label{sec:architecture}
The automated pipeline processes incoming high-resolution images through a series of specialized stages. The initial stage handles structural localization to isolate key regions of interest from the background clutter of the wagon body.

\begin{figure}[htbp]
\centering
\imageplaceholder{0.45\textwidth}{60pt}
\caption{System architecture flow diagram showing localization, OCR routing, and database verification.}
\label{fig:flowchart}
\end{figure}

\subsection{Object Detection Framework (YOLOv11m)}
We chose the YOLOv11m architecture for parameter localization and safety stencil classification due to its balance of inference speed and bounding box accuracy. YOLOv11m introduces several architectural enhancements over previous iterations:
\begin{itemize}
    \item \textbf{C3k2 Bottleneck Block:} This block improves gradient flow and representation capacity by incorporating Cross Stage Partial bottlenecks with variable kernel configurations.
    \item \textbf{SPPF (Spatial Pyramid Pooling - Fast):} This sub-module encodes multi-scale features, capturing global contextual features with minimal computational overhead.
    \item \textbf{Decoupled Head:} The network uses separate branches for classification and bounding box regression (Fig. \ref{fig:decoupled_head}), enabling more focused feature learning.
\end{itemize}

\begin{figure}[htbp]
\centering
\imageplaceholder{0.45\textwidth}{40pt}
\caption{Structure of the decoupled detection head in YOLOv11m separating classification, regression, and confidence branches.}
\label{fig:decoupled_head}
\end{figure}

\subsection{Parameter Localization Configuration}
The localization model is trained to detect five specific parameter classes:
\begin{enumerate}
    \item \textbf{Wagon Number:} The 11-digit identification mark unique to each wagon.
    \item \textbf{Wagon Type:} The classification code (e.g., BOXN, BOBYN, BCN) representing the wagon structure.
    \item \textbf{Owner Name:} The owning railway administration (e.g., IR for Indian Railways, CONCOR, or private companies).
    \item \textbf{Zone Marking:} The zonal railway abbreviation (e.g., SEC, SE, SER, ECR).
    \item \textbf{Tare Load Details:} The empty weight of the wagon, which is crucial for load calculations.
\end{enumerate}

\subsection{Safety Stencil Classification}
Wagons must undergo periodic maintenance (POH - Periodic Overhaul, and ROH - Routine Overhaul). The dates and locations of these inspections are stenciled onto the wagon body. A secondary, custom-trained YOLOv11m classification model scans the wagon crops to audit the presence of these stencils, flagging wagons with missing or illegible maintenance markers.

\section{Multi-Tier OCR Ensemble Strategy}
\label{sec:ocr_ensemble}
Text printed on wagon steel panels is frequently degraded. To prevent recognition failures, we developed a 5-tier fallback OCR ensemble strategy (summarized in Table \ref{tab:ocr_performance_comparison}). The system sequentially routes each localized text crop through these tiers until a valid string is confirmed.

\subsection{Tier 1: Preprocessed Crop with PaddleOCR}
The first tier applies image processing filters to enhance contrast before character extraction:
\begin{enumerate}
    \item \textbf{Radon Transform Deskewing:} The rotation skew angle $\theta$ of the text crop is calculated, and the image is rotated to align the text horizontally.
    \item \textbf{LAB CLAHE Enhancement:} The crop is converted from the BGR color space to the LAB color space. Contrast Limited Adaptive Histogram Equalization (CLAHE) is applied to the L-channel (luminance) to boost contrast without amplifying background metal grain noise:
    \begin{equation}
    g_{i,j} = \left[ g_{max} - g_{min} \right] \cdot P(f_{i,j}) + g_{min}
    \end{equation}
    Where $P(f_{i,j})$ is the cumulative distribution function of the local neighborhood histogram, clipped to limit local contrast stretching.
    \item \textbf{Bilinear Rescaling:} The crop is rescaled to a target height of 80 pixels to provide optimal character resolution.
\end{enumerate}

\subsection{Tier 2: Raw Crop with PaddleOCR}
If the preprocessing filters over-segment the characters (which can happen with thin, weathered lettering), Tier 2 processes the raw crop. This layer includes adaptive margins, expanding the YOLO bounding box boundaries by 5\% to prevent edge clipping.

\subsection{Tier 3: EasyOCR Fallback}
If PaddleOCR fails, the crop is passed to EasyOCR. EasyOCR combines a ResNet backbone with a Bidirectional LSTM (Long Short-Term Memory) network to capture horizontal character context, making it effective at resolving broken letters on weathered paint.

\subsection{Tier 4: Tesseract OCR with Whitelists}
Tier 4 utilizes Tesseract OCR. We configure it in Page Segmentation Mode 7 (PSM 7), which treats the crop as a single text line. We restrict character recognition using strict alphanumeric whitelists:
\begin{itemize}
    \item \textbf{Wagon Number whitelist:} \texttt{0-9} (numeric characters only).
    \item \textbf{Zone / Owner whitelist:} \texttt{A-Z} (uppercase alphabetic characters only).
\end{itemize}

\subsection{Tier 5: Full Image Fallback Scan}
If the localization model fails to detect the wagon number bounding box due to perspective angles or dust coverage, Tier 5 performs a sliding-window text line scan across the full high-resolution image to locate the 11-digit pattern.

\begin{table*}[t]
\caption{Comparison of OCR Fallback Tiers on Degraded Metal Markings}
\begin{center}
\begin{tabular}{lllll}
\toprule
\textbf{Tier} & \textbf{Core Engine} & \textbf{Processing Applied} & \textbf{Target Error Mode Resolved} & \textbf{Relative Latency} \\
\midrule
Tier 1 & PaddleOCR (PP-OCRv4) & Deskew + LAB CLAHE + Resizing & Low local contrast, skew, surface rust & Baseline (1.0x) \\
Tier 2 & PaddleOCR (PP-OCRv4) & Raw crop + 5\% adaptive margins & Character clipping at bounding box boundaries & 1.05x \\
Tier 3 & EasyOCR & ResNet + BiLSTM + CTC Decoder & Weathered paint peeling, split characters & 1.85x \\
Tier 4 & Tesseract & PSM 7 + Strict Whitelisting & Digit/character substitution errors & 0.65x \\
Tier 5 & PaddleOCR (PP-OCRv4) & Full-frame sliding window scan & Bounding box localization failure & 4.20x \\
\bottomrule
\end{tabular}
\label{tab:ocr_performance_comparison}
\end{center}
\end{table>

\section{S4 Mod-10 Checksum and Candidate Correction}
\label{sec:checksum}
The 11-digit wagon number contains built-in redundancy to verify correctness. We exploit this using the S4 Mod-10 checksum algorithm.

\subsection{Mathematical Formulation}
Let $D$ be the 11-digit wagon number string:
\begin{equation}
D = d_1 d_2 d_3 d_4 d_5 d_6 d_7 d_8 d_9 d_{10} d_{11}
\end{equation}
Where $d_i \in \{0, 1, \dots, 9\}$. The digits are weighted based on their position:
\begin{equation}
w_i = 
\begin{cases} 
3 & \text{if } i \text{ is odd} \\
1 & \text{if } i \text{ is even}
\end{cases}
\end{equation}

The validation sum $S$ is computed as:
\begin{equation}
S = \sum_{i=1}^{10} w_i \cdot d_i
\end{equation}
Or, written out:
\begin{equation}
S = 3(d_1 + d_3 + d_5 + d_7 + d_9) + 1(d_2 + d_4 + d_6 + d_8 + d_{10})
\end{equation}

Let the modulo 10 remainder of $S$ be:
\begin{equation}
R = S \pmod{10}
\end{equation}

The expected check digit $C$ is defined by:
\begin{equation}
C = (10 - R) \pmod{10}
\end{equation}

The wagon number is mathematically valid if and only if:
\begin{equation}
d_{11} = C
\end{equation}

\subsection{Combinatorial Candidate Correction}
If the calculated check digit does not match $d_{11}$, the system flags the extraction as failed. Instead of discarding the result, the correction engine generates candidate replacements.

\begin{figure}[htbp]
\centering
\imageplaceholder{0.45\textwidth}{50pt}
\caption{Fuzzy matching candidate lookup logic using Levenshtein distance against the ROAMS zone-owner matrix.}
\label{fig:fuzzy_logic}
\end{figure}

The algorithm identifies common OCR substitutions (e.g., replacing '8' with 'B' or '6' with 'G') and builds a candidate search space. It calculates the Levenshtein distance between the extracted candidate prefixes and the valid configurations in the ROAMS database:
\begin{equation}
L(u, v) = \min \left( \text{cost of transforming } u \text{ to } v \right)
\end{equation}

By checking candidate numbers against the valid railway type and owner matrices (e.g., verifying if the first two digits match the wagon type, and the next two digits match the zone/owner), the correction engine resolves digit misreads and recovers the correct number.

\section{System Deployment and Stability}
The system is built as a distributed microservice architecture to ensure high throughput and prevent thread blocks.

\subsection{Decoupled GPU OCR Microservice}
Running heavy text recognition libraries (like PyTorch and PaddlePaddle) directly inside the FastAPI web server thread pool can cause CUDA memory locks and slow down request handling. 

To solve this, we decoupled the OCR pipeline. The main FastAPI server runs on port 8000, while the heavy OCR models are hosted on an isolated microservice running on port 8001. All text extraction requests are sent to this service via HTTP POST calls.

\subsection{oneDNN Thread Safety on Windows}
During local execution on Windows, PaddlePaddle 3.x's static graph execution engine can trigger memory crashes in oneDNN (MKLDNN) libraries. To prevent these crashes, we configure the OCR initialization with:
\begin{equation}
\text{enable\_mkldnn} = \text{False}
\end{equation}
And set the following environment variables at startup:
\begin{itemize}
    \item \texttt{FLAGS\_enable\_pir\_api=0}
    \item \texttt{FLAGS\_use\_mkldnn=0}
    \item \texttt{FLAGS\_enable\_pir\_in\_executor=0}
\end{itemize}
This ensures stable, crash-free execution on Windows.

\subsection{React Dashboard Frontend}
The React web application communicates with the FastAPI backend to visualize auditing runs. It includes:
\begin{enumerate}
    \item \textbf{Map Pane:} Interactive map showing geotagged wagon locations.
    \item \textbf{OCR Process Pane:} Displays crop steps, showing the processing stages and confidence scores for each detected parameter.
    \item \textbf{Wagon Manifest Table:} A tabular log displaying the parameters, validation flags, and stencil statuses.
\end{enumerate}

\section{Results and Performance Benchmarking}
We evaluated the system on an evaluation dataset consisting of easy, medium, and hard wagon images. Hard cases included weathered surfaces, heavy rust, shadow lines, and perspective skew.

\begin{table}[htbp]
\caption{Extraction Accuracy by Parameter Class}
\begin{center}
\begin{tabular}{lccc}
\toprule
\textbf{Parameter Class} & \textbf{Easy Cases} & \textbf{Medium Cases} & \textbf{Hard Cases} \\
\midrule
Wagon Number   & 100.0\% & 96.0\% & 86.0\% \\
Wagon Type     & 100.0\% & 98.0\% & 91.0\% \\
Owner Name     & 98.5\%  & 94.0\% & 88.0\% \\
Zone Marking   & 98.5\%  & 95.0\% & 89.0\% \\
Tare Details   & 96.0\%  & 90.0\% & 80.0\% \\
Safety Stencil & 97.0\%  & 92.0\% & 85.0\% \\
\bottomrule
\end{tabular}
\label{tab:class_accuracies}
\end{center}
\end{table}

Table \ref{tab:class_accuracies} displays the extraction accuracies across different cases. While easy images achieve near-perfect scores, hard images benefit from the OCR fallback tiers and correction logic, maintaining an accuracy of 86.0\% for the 11-digit wagon number.

\subsection{Inference Speed Benchmarks}
Inference benchmarks were conducted on an NVIDIA A4500 GPU. The processing time breakdown per frame is detailed in Table \ref{tab:speed_benchmarks}.

\begin{table}[htbp]
\caption{Latency Breakdown (NVIDIA A4500 GPU)}
\begin{center}
\begin{tabular}{lc}
\toprule
\textbf{Pipeline Stage} & \textbf{Average Execution Time} \\
\midrule
YOLOv11m Parameter Localization & 42 ms \\
YOLOv11m Stencil Auditing       & 38 ms \\
Tier 1 Preprocessing + OCR      & 420 ms \\
Tier 2 Fallback OCR (If run)    & 410 ms \\
Tier 3 EasyOCR Fallback (If run)& 860 ms \\
S4 Checksum & 2 ms \\
Combinatorial Correction       & 12 ms \\
\bottomrule
\end{tabular}
\label{tab:speed_benchmarks}
\end{center}
\end{table}

The average end-to-end processing time for cases resolved at Tier 1 is approximately 514 ms, enabling real-time audits.

\section{Conclusion}
This paper presented a multi-tiered deep learning pipeline for automated wagon parameter auditing and safety compliance checks for Indian Railways. By combining custom YOLOv11m object detectors with a 5-tier OCR ensemble and S4 Mod-10 candidate correction, the system achieves robust extraction performance under challenging yard conditions. Decoupling the backend into a FastAPI service and an isolated GPU text-recognition microservice prevents CUDA memory conflicts and ensures execution stability. Future work will investigate model quantization to enable edge deployment on low-power devices.

\begin{thebibliography}{00}
\bibitem{yolo11} Ultralytics, "YOLO11 Object Detection and Segmentation Models," 2024. [Online]. Available: https://github.com/ultralytics/ultralytics.
\bibitem{paddleocr} Y. Du et al., "PP-OCR: A Practical Ultra Lightweight OCR System," \textit{arXiv preprint arXiv:2009.09941}, 2020.
\bibitem{easyocr} JaidedAI, "EasyOCR: Ready-to-use OCR with 80+ Languages Supported," 2023. [Online]. Available: https://github.com/JaidedAI/EasyOCR.
\bibitem{levenshtein} V. I. Levenshtein, "Binary codes capable of correcting deletions, insertions, and reversals," \textit{Soviet Physics Doklady}, vol. 10, no. 8, pp. 707-710, 1966.
\end{thebibliography}

\end{document}
